\section*{Chapitre \ref{ch:b1:ecosystem-organization} --- Écosystèmes et structure trophique}

\begin{solution}{exo:b1:ecosystem-organization:1}
L'\omterm{def:b1:ecosystem-organization:ecosystem}{écosystème} de la mare, c'est la \omterm{def:b1:ecosystem-organization:ecosystem}{biocénose} --- toutes ses \omterm{def:b1:populations:population}{populations}~:
algues, potamots, escargots, larves d'insectes, grenouilles, poissons,
bactéries --- avec le \omterm{def:b1:ecosystem-organization:ecosystem}{biotope}~: l'eau, sa température, la lumière, les
gaz dissous et les minéraux, la vase.
\end{solution}

\begin{solution}{exo:b1:ecosystem-organization:2}
\omterm{def:b1:ecosystem-organization:niche}{Habitat}~: le lieu où vit une espèce (une pessière). \omterm{def:b1:ecosystem-organization:niche}{Niche}~: sa manière
d'y vivre --- ce qu'elle mange, quand, où dans l'arbre, ce qu'elle
tolère. \omterm{def:b1:ecosystem-organization:niche}{Niche} fondamentale~: l'intervalle qu'elle pourrait occuper
seule~; réalisée~: ce que lui laissent concurrents et prédateurs.
\end{solution}

\begin{solution}{exo:b1:ecosystem-organization:3}
$20\,800/1.7\times 10^6 = \qty{1.2}{\%}$~; les herbivores ont prélevé
$3400/8800 = \qty{39}{\%}$ de la production nette.
\end{solution}

\begin{solution}{exo:b1:ecosystem-organization:4}
\omterm{def:b1:ecosystem-organization:trophic}{Producteurs}, \omterm{def:b1:ecosystem-organization:trophic}{consommateurs} (primaires, secondaires, tertiaires),
\omterm{def:b1:ecosystem-organization:trophic}{décomposeurs}, détritivores~; ce sont les \omterm{def:b1:ecosystem-organization:trophic}{décomposeurs} qui rendent aux
\omterm{def:b1:ecosystem-organization:trophic}{producteurs} les minéraux de la matière morte.
\end{solution}

\begin{solution}{exo:b1:ecosystem-organization:5}
PPN $= 800\times 17 = \qty{13.6}{MJ/m^2}$~; mangé $\qty{2.04}{MJ}$~;
assimilé $\qty{1.02}{MJ}$~; lapin $\qty{51}{kJ/m^2}$, soit
\qty{510}{MJ} par hectare (\qty{30}{kg} de lapin à \qty{17}{kJ/g}).
Efficacité $51/13\,600 = \qty{0.4}{\%}$.
\end{solution}

\begin{solution}{exo:b1:ecosystem-organization:6}
La biomasse est un stock~: un phytoplancton mangé aussi vite qu'il
pousse n'a jamais besoin de s'accumuler, si bien que le stock présent
de \omterm{def:b1:ecosystem-organization:trophic}{producteurs} peut être inférieur à celui des animaux plus
longévifs qui en vivent. L'énergie est un flux~: chaque niveau reçoit du
niveau inférieur tout ce qu'il aura jamais et doit en perdre une part
en chaleur, si bien que le flux ne peut que décroître vers le haut.
\end{solution}

\begin{solution}{exo:b1:ecosystem-organization:7}
$p = 0.5, 0.3, 0.15, 0.05$~: $H = -(0.5\ln 0.5 + 0.3\ln 0.3 + 0.15\ln
0.15 + 0.05\ln 0.05) = 0.347 + 0.361 + 0.285 + 0.150 = 1.14$~;
équitabilité $1.14/\ln 4 = 0.82$. Quatre espèces à 25 chacune~: $H = \ln 4 = 1.39$,
équitabilité 1.
\end{solution}

\begin{solution}{exo:b1:ecosystem-organization:8}
Nette $= \qty{3}{mg}$ d'$\mathrm{O_2}$ par litre et par jour~;
respiration \qty{1}{mg}~; brute $3 + 1 = \qty{4}{mg}$. En carbone~: brute
\qty{1.5}{mg}, nette \qty{1.1}{mg} de C par litre et par jour.
\end{solution}

\begin{solution}{exo:b1:ecosystem-organization:9}
Chaque maillon n'en garde qu'un dixième~: dix maillons ne laisseraient
qu'un milliardième de la production primaire, bien trop peu pour
nourrir un seul animal sur quelque surface que ce soit. Là où la
production est faible (désert) ou l'aire petite (île), le
dixième-du-dixième s'épuise après moins de maillons~; la forte
production d'une forêt en soutient un ou deux de plus.
\end{solution}

\begin{solution}{exo:b1:ecosystem-organization:10}
Grain~: \qty{1}{t} demande un sixième d'hectare. Bœuf~: \qty{1}{t}
d'énergie sous forme de bœuf demande \qty{10}{t} de grain, soit
\num{1.7} hectare --- dix fois la surface. Manger au second niveau coûte
l'équivalent d'un niveau d'énergie~; une planète nourrit dix fois plus
de gens avec du grain qu'avec de la viande nourrie au grain, et le
\omterm{def:b1:ecosystem-organization:trophic}{niveau trophique} humain, moyenné sur les régimes, pèse lourd dans la
\omterm{thm:b1:populations:logistic}{capacité de charge}.
\end{solution}

\begin{solution}{exo:b1:ecosystem-organization:11}
Par mètre carré, l'océan est un désert, mais il couvre les deux tiers
de la planète~: un faible taux multiplié par une aire énorme donne la
moitié du total. Sa production est limitée par les éléments nutritifs,
qui coulent hors de la couche éclairée~; là où les courants les
ramènent (upwellings côtiers, plateaux), la production est dix fois plus
forte, et c'est là, sur quelques pour cent de la mer, que se
concentrent les poissons et les pêcheries.
\end{solution}

\begin{solution}{exo:b1:ecosystem-organization:12}
L'énergie entre sous forme de lumière, traverse quelques niveaux
d'\omterm{def:b1:organism-environment:organism}{organismes} et repart entièrement en chaleur~: en ce sens,
l'\omterm{def:b1:ecosystem-organization:ecosystem}{écosystème} convertit bien la lumière solaire en chaleur, et les
\omterm{def:b1:ecosystem-organization:trophic}{décomposeurs}, qui respirent tout ce que les \omterm{def:b1:ecosystem-organization:trophic}{consommateurs} ont laissé,
sont l'endroit où la plus grande part s'en va. Mais la matière n'est
pas consommée~: carbone, azote et phosphore circulent en cycle, rendus
par ces mêmes \omterm{def:b1:ecosystem-organization:trophic}{décomposeurs} aux \omterm{def:b1:ecosystem-organization:trophic}{producteurs}, si bien que le système
persiste~; et les \omterm{def:b1:organism-environment:organism}{organismes} ne sont pas un sous-produit mais la
structure qui organise le flux --- les \omterm{def:b1:ecosystem-organization:niche}{niches}, les réseaux et les
pyramides sont ce que l'énergie bâtit en chemin. La phrase a la
thermodynamique et manque la biologie.
\end{solution}

\begin{solution}{pb:b1:ecosystem-organization:1}
\textbf{1.} $20\,800 - 12\,000 = \qty{8800}{kcal}$ par mètre carré et
par an.
\textbf{2.} \qty{1.2}{\%}. Les pertes~: lumière non absorbée (la moitié
du spectre, réflexion, eau), \omterm{def:b1:flowering-plant-organization:organs}{feuilles} saturées en plein soleil ou à
l'ombre, photorespiration, et la respiration des plantes elles-mêmes
(déjà comptée entre le brut et le net).
\textbf{3.} $12\,000/20\,800 = \qty{58}{\%}$.
\textbf{4.} Mangé \qty{3400}{}~; directement aux \omterm{def:b1:ecosystem-organization:trophic}{décomposeurs}
$8800 - 3400 = \qty{5400}{kcal}$.
\textbf{5.} $8800/10 = \qty{880}{g/m^2}$~: comme une forêt tempérée ou
une prairie riche.
\textbf{6.} $1.7\times 10^6 - 20\,800 = \qty{1.68e6}{kcal}$, soit
\qty{98.8}{\%} du total~: réfléchie, transmise, ou absorbée par l'eau et
la roche puis réémise en chaleur.
\textbf{7.} Production $1500 - 1100 = \qty{400}{kcal}$~; efficacité
d'assimilation $1500/3400 = \qty{44}{\%}$.
\textbf{8.} $400/1500 = \qty{27}{\%}$.
\textbf{9.} $400/8800 = \qty{4.5}{\%}$ (\qty{16}{\%} si l'on compte sur
ce qui a été mangé, \qty{3400}{}).
\textbf{10.} Production des carnivores $380 - 300 = \qty{80}{kcal}$~;
efficacité $80/400 = \qty{20}{\%}$ (ils ont mangé \qty{380}{} des
\qty{400}{} produites).
\textbf{11.} $6/80 = \qty{7.5}{\%}$.
\textbf{12.} \qty{6}{kcal}~; $6/1.7\times 10^6 = \num{3.5e-6}$, quelques
millionièmes.
\textbf{13.} La viande est digestible et proche du mangeur par sa
composition~; la matière végétale est fibreuse et largement
indigestible, si bien qu'un herbivore perd une plus grande part de son
ingestion en fèces.
\textbf{14.} Production nette non mangée \qty{5400}{}~; fèces des
herbivores $3400 - 1500 = \qty{1900}{}$~; production d'herbivores non
mangée $400 - 380 = \qty{20}{}$~; fèces des carnivores (l'assimilation
n'est pas donnée~; prendre mangé moins respiré moins production $= 0$,
tout assimilé) et production de carnivores non mangée
$80 - 20 = \qty{60}{}$~; production des carnivores supérieurs
\qty{6}{}~: total \qty{7386}{kcal}.
\textbf{15.} $12\,000 + 1100 + 300 + 14 + 7386 = \qty{20800}{kcal}$.
\textbf{16.} Égale à la production brute~: la source est en régime
stationnaire et ne stocke aucune biomasse d'une année sur l'autre.
\textbf{17.} $7386/20\,800 = \qty{36}{\%}$ (l'essentiel du reste est la
respiration des \omterm{def:b1:ecosystem-organization:trophic}{producteurs} eux-mêmes).
\textbf{18.} \omterm{def:b1:ecosystem-organization:trophic}{Producteurs} \qty{8800}{}, herbivores \qty{400}{},
carnivores \qty{80}{}, carnivores supérieurs \qty{6}{}~; efficacités
\qty{4.5}{\%}, \qty{20}{\%}, \qty{7.5}{\%}.
\textbf{19.} \qty{8000}{}, \qty{400}{}, \qty{100}{}, \qty{15}{kcal/m^2}~:
droite, chaque niveau valant un dixième ou moins du précédent.
\textbf{20.} \omterm{def:b1:ecosystem-organization:trophic}{Producteurs} $8000/8800 = 0.9$ an~; herbivores $400/400 =
1$ an~; carnivores $100/80 = 1.25$ an~; carnivores supérieurs $15/6 =
2.5$ ans. Les plantes se renouvellent le plus vite~: \omterm{def:b1:body-plans-tissues:tissue}{tissus} de courte
durée, broutés en continu~; les carnivores supérieurs sont des animaux
longévifs dont le stock se renouvelle lentement.
\textbf{21.} Un héron de \qty{2}{kg} pèse environ \qty{500}{g} sec,
soit \qty{5000}{kcal}, et il lui faut à peu près autant de production
chaque année pour se remplacer, lui et ses jeunes~:
$5000/6 \approx \qty{800}{m^2}$ de source --- en pratique bien
davantage, puisqu'il respire aussi~: à \qty{100}{kcal} par jour,
\qty{36500}{kcal} d'ingestion par an, c'est-à-dire \qty{6000}{m^2} du
niveau inférieur (production des carnivores \qty{80}{}), ce qui
ressemble bien au territoire d'un héron.
\textbf{22.} La production nette supplémentaire n'est pas mangée et va
aux \omterm{def:b1:ecosystem-organization:trophic}{décomposeurs}, dont la respiration double~; la nuit, sans
photosynthèse, leur demande en oxygène peut vider l'eau de son oxygène
et tuer les poissons --- l'eutrophisation.
\textbf{23.} La production des herbivores tombe à \qty{200}{}~: les
carnivores reçoivent la moitié de leur nourriture, leur production
tombe vers \qty{40}{}, celle des carnivores supérieurs vers 3, et il y
a moins de hérons~; les zostères, moins broutées, s'épaississent et
passent en plus grande part, non mangées, aux \omterm{def:b1:ecosystem-organization:trophic}{décomposeurs}.
\textbf{24.} L'énergie des carnivores supérieurs est minuscule, mais
leur effet porte sur les effectifs des carnivores, qui contrôlent les
herbivores, qui contrôlent l'herbe~: les retirer laisse monter les
carnivores, baisser les herbivores et croître l'herbe --- une cascade à
travers le réseau (\cref{ch:b1:species-interactions})~; retirer
quelques kilocalories d'herbivores ne change rien que quelques
kilocalories.
\textbf{25.} Des \omterm{def:b1:ecosystem-organization:trophic}{producteurs} aux herbivores \qty{4.5}{\%} (de la PPN~;
\qty{16}{\%} de ce qui a été mangé), des herbivores aux carnivores
\qty{20}{\%}, des carnivores aux carnivores supérieurs \qty{7.5}{\%}~;
de la lumière solaire à la production brute \qty{1.2}{\%}~; quelques
millionièmes de la lumière solaire finissent dans les carnivores
supérieurs.
\end{solution}
